\section{Introduction}
Economic and financial time series often exhibit persistent trends alongside cyclical dynamics whose amplitude, duration, and frequency evolve over time. Structural change, policy interventions, and market regime shifts can produce behavior that is difficult to capture with models that assume stationary dynamics or a fixed spectral structure.
Structural time-series and state-space approaches address part of this problem by decomposing observations into latent components such as trend and cycle \cite{harvey1989structural,durbin2012statespace}. However, these models often impose fixed or weakly time-varying cycle frequencies. Time-varying parameter models and regime-switching frameworks partially relax these assumptions, but typically remain linear and rely on strong parametric structure \cite{lubik2015tvpvar,hamilton1989regime}.
Recent work on implicit neural representations shows that sinusoidal nonlinearities can represent complex oscillatory patterns efficiently \cite{sitzmann2020siren}, while the Hilbert--Huang Transform describes nonstationary signals through time-varying amplitude and instantaneous frequency \cite{huang1998hht}. Motivated by these ideas, we introduce the \textit{Non-Stationary Spectral Decomposition Network} (NS-SDN), a neural state-space architecture that models a time series as a sum of sinusoidal components with amplitude, frequency, and phase driven by a latent dynamical state.
